\exercise{Counting spanning trees}{2}
Use the matrix-tree theorem to count the number of spanning trees in a fully connected network of 4 nodes. Can you also count them in a fully connected network of $n$ nodes?

\solution 
The Laplacian matrix of a fully-connected network of four nodes is 
\eq{
{\bf L_4} = \avecccc{
 3 &-1 &-1 &-1 \\
-1 & 3 &-1 &-1 \\
-1 &-1 & 3 &-1 \\
-1 &-1 &-1 & 3 }
}
So one way of counting the spanning trees are to compute the eigenvalues. Building on Chap.~20 we can guess the eigenvectors 
\eq{
\avecccc{
 3 &-1 &-1 &-1 \\
-1 & 3 &-1 &-1 \\
-1 &-1 & 3 &-1 \\
-1 &-1 &-1 & 3 }
\avec{1\\ 1\\ 1\\ 1} = \avec{0 \\ 0 \\ 0 \\ 0}
}
\eq{
\avecccc{
 3 &-1 &-1 &-1 \\
-1 & 3 &-1 &-1 \\
-1 &-1 & 3 &-1 \\
-1 &-1 &-1 & 3 }
\avec{1\\ -1\\ 0\\ 0} = \avec{4 \\ -4 \\ 0 \\ 0}
}
\eq{
\avecccc{
 3 &-1 &-1 &-1 \\
-1 & 3 &-1 &-1 \\
-1 &-1 & 3 &-1 \\
-1 &-1 &-1 & 3 }
\avec{0\\ 1\\ -1\\ 0} = \avec{0 \\ 4 \\ -4 \\ 0}
}
\eq{
\avecccc{
 3 &-1 &-1 &-1 \\
-1 & 3 &-1 &-1 \\
-1 &-1 & 3 &-1 \\
-1 &-1 &-1 & 3 }
\avec{0\\ 0\\ 1\\ -1} = \avec{0 \\ 0 \\ 4 \\ -4}
}
Hence the eigenvalues are 
\eq{
\lambda_1=0,\quad \lambda_2=4,\quad \lambda_3=4, \quad \lambda_4=4.
}
So we can count the number of spanning trees using
\eq{
n_{\rm Trees} = \frac{\lambda_2\lambda_3\lambda_4}{n} = \frac{4\cdot 4 \cdot 4}{4} = 16
}
which is correct, but we wanted to get the result using the matrix-tree theorem. So let's try another way. 

Removing one row and the corresponding column we are left with 
\eq{
{\bf L_{4,\overline{11}}} = \aveccc{
 3 &-1 &-1 \\
-1 & 3 &-1 \\
-1 &-1 & 3 }
}
We can now compute 
\eqa{
n_{\rm Trees} &=& |{\bf L_{4,\overline{11}}}| \\
  &=& 27 - 1 - 1 - 3 - 3 - 3 \\
  &=& 16
}
which again gives us the result. But to generalize to larger systems it is better to move to eigenvalues rather then computing the determinant directly. Using the symmetry we can guess the eigenvectors 
\eq{
\vec{v_1} = \avec{1\\ 1\\ 1}, \quad \vec{v_2}=\avec{1\\ -1\\ 0},\quad 
\vec{v_3}=\avec{0\\ 1 \\ -1},
}
Multiplying these with ${\bf L_{4,\overline{11}}}$ confirms that these are indeed eigenvectors and reveals the corresponding eigenvalues 
\eq{
\lambda_1 = 1, \lambda_2=4, \lambda_3=4 
}
Hence the number of spanning trees is 
\eq{
n_{\rm s.t.} = 1\cdot 4 \cdot 4 = 16
}
If you want to verify this result in yet another way you can do it as follows: As there are four nodes, every spanning tree will have three links. Three links can eiher form a start or chain. On the fully connected network of four nodes, we can construct 4 different stars (one with each of the four different nodes in the center. To count the number of chains with three links, consider that we have four choices for the starting node, then three choices for the second node in the chain, then two choices for the third node and then there is only one choice left for the final nodes, so $4\cdot 3\cdot 2 \cdot 1 = 24$, but now we have counted each chain in both directions, so we divide by 2 to correct for the overcounting which gives us a total of 12 different chains. If we add the 12 spanning chains to the 4 spanning stars, we get 16 spanning trees in total.)

In the general case of a fully connected network of $n$ nodes, the reduced matrix is a $(n-1)\times (n-1)$ matrix which has one eigenvector of the form of $\vec{v_1}$ with an eigenvalues of 1, and $n-2$ eigenvectors of the form of $\vec{v_2}$ and $\vec{v_3}$ with an eigenvalue of $n$. 

Hence the number of spanning trees is 
\eq{
n_{\rm s.t.} = 1 \cdot \overbrace{n \cdot n \cdot \ldots}^{\mbox{$n-2$ terms}} = n^{n-2}. 
}
This is Cayley's Tree formula. 

